\documentclass[10pt]{article}
\usepackage[utf8]{inputenc}
\usepackage[T1]{fontenc}
\usepackage{graphicx}
\usepackage[export]{adjustbox}
\graphicspath{ {./images/} }
\usepackage{amsmath}
\usepackage{amsfonts}
\usepackage{amssymb}
\usepackage[version=4]{mhchem}
\usepackage{stmaryrd}
\usepackage{bbold}
\usepackage{caption}

\begin{document}
\captionsetup{singlelinecheck=false}
Uniqueness, Determinacy, Genericity, Stability

\section*{Uniqueness, Determinacy, Genericity}
\begin{itemize}
  \item Comparative statics (CS): How does a WE change when the underlying parameters of the economy change?
  \item Meaningful comparative statics in two cases,
  \item Uniqueness of equilibria and
  \item Local uniqueness plus a nice transition.\\
\includegraphics[max width=\textwidth, alt={}, center]{5e8d08d5-9594-42e7-ba77-0f861991fae4-04_857_780_181_145}\\
\includegraphics[max width=\textwidth, alt={}, center]{5e8d08d5-9594-42e7-ba77-0f861991fae4-04_841_768_185_973}\\
\includegraphics[max width=\textwidth, alt={}, center]{5e8d08d5-9594-42e7-ba77-0f861991fae4-04_853_776_185_1797}
\end{itemize}

\section*{1. Uniqueness of WE and comparative statics}
\begin{itemize}
  \item If WE is unique (before and after a change in parameters), the policy change is univocally associated to a given change in the Walrasian equilibrium.
  \item There are no uniqueness results under reasonable assumptions on individual preferences.
  \item See Mas-Colell, Whinston, and Green (1995), 17.F.
\end{itemize}

\section*{2. Local uniqueness and comparative statics}
Let $E q(\mathcal{E})$ be the set of equilibrium prices of the economy $\mathcal{E}$.\\
2.a. Local uniqueness of WE: $\forall p^{*} \in E q(\mathcal{E}), \exists \delta>0$ such that there is no equilibrium price $q^{*} \in \Delta$ with $q^{*} \neq p^{*},\left\|q^{*}-p^{*}\right\|<\delta$ and\\
2.b. For sufficiently small change in parameters, the number of equilibria is unchanged and each equilibrium moves nicely (i.e. continuously or differentiably) as the parameter changes.

These requirements are not sufficient for a meaningul comparative static exercise.

\section*{Main findings}
\begin{enumerate}
  \item Identification of a class of economies (called regular) where 2.a and 2.b hold (i.e. there are finitely many WE and the WE vary nicely with the parameters of the economy.
  \item Genericity: "A large number of economies are within the identified class" Debreu (1970) using Sard's Theorem.
  \item Link to stability of equilibria. Index Theorem (Poincar'e-Hopf Theorem).
\end{enumerate}

\section*{Preliminaries}
Definition $\mathbf{1}$ Let $\mathcal{E}$ be a pure exchange economy. Agent $i$ 's excess demand is $z: \mathbb{R}_{++}^{L} \times \mathbb{R}^{L}$,

$$
z_{i}\left(p, \omega_{i}\right)=D_{i}\left(p, \omega_{i}\right)-\omega_{i}
$$

The aggregate excess demand is $z: \mathbb{R}_{++}^{L} \times \mathbb{R}^{L I}$,

$$
\begin{aligned}
z(p, \omega) & =\sum_{i=1}^{I} z_{i}\left(p, \omega_{i}\right)=\sum_{i=1}^{I} D_{i}\left(p, \omega_{i}\right)-\bar{\omega} \\
z_{\ell}(p, \omega) & =\sum_{i=1}^{I} z_{i \ell}\left(p, \omega_{i}\right)
\end{aligned}
$$

\section*{Local uniqueness and CS, example}
$L=2$; Walras' law; $\left(\forall p \in \mathbb{R}_{++}^{2}\right) \hat{p}=\frac{p_{1}}{p_{2}}$ and $\hat{z}(\hat{p}, \omega)=z_{1}\left(\frac{p_{1}}{p_{2}}, 1, \omega\right)$.\\
\includegraphics[max width=\textwidth, alt={}, center]{5e8d08d5-9594-42e7-ba77-0f861991fae4-09_1410_2543_512_140}

\section*{Local uniqueness but}
\begin{center}
\includegraphics[max width=\textwidth, alt={}]{5e8d08d5-9594-42e7-ba77-0f861991fae4-10_1396_2536_345_143}
\end{center}

\begin{figure}[h]
\begin{center}
\captionsetup{labelformat=empty}
\caption{No local uniqueness}
  \includegraphics[alt={},max width=\textwidth]{5e8d08d5-9594-42e7-ba77-0f861991fae4-11_1399_2536_342_143}
\end{center}
\end{figure}

\section*{Local uniqueness-Maintained assumptions}
\begin{itemize}
  \item Let $\mathcal{E}$ be a pure exchange economy and $z$ its aggregate excess demand.
  \item $z$ satisfies the assumptions of the DGKN Lemma.
  \item $z$ is homogeneous of degree zero,
\end{itemize}

$$
\left(\forall p \in \mathbb{R}_{++}^{L}, \forall \alpha>0\right) \quad z(\alpha p)=z(p)
$$

\begin{itemize}
  \item $z \in C^{1}$.
\end{itemize}

\section*{Preliminaries}
\begin{itemize}
  \item Normalize prices so that $p_{L}=1$, let $\hat{p}=\left(p_{1}, \ldots, p_{L-1}\right) \in \mathbb{R}_{++}^{L-1}$.
  \item Let $\hat{z}: \mathbb{R}_{++}^{L-1} \rightarrow \mathbb{R}^{L-1}$ be $\hat{z}(\hat{p})=\left(z_{1}(\hat{p}, 1), \ldots, z_{L-1}(\hat{p}, 1)\right)$.
  \item \hspace{0pt} [Walras' law and $p \gg 0] \Longrightarrow(\hat{z}(\hat{p})=0 \Longleftrightarrow z(\hat{p}, 1)=0)$
  \item The Jacobian matrix of $\hat{z}, D \hat{z}$ is $L-1 \times L-1$.
\end{itemize}

\section*{Regularity}
Definition 2 (Regular) An equilibrium price $p^{*}$ of the economy $\mathcal{E}$ is regular if $\left.\operatorname{det} D \hat{z}\right|_{\hat{p}^{*}} \neq 0 . \backslash$\\
The economy $\mathcal{E}$ is regular if all its equilibrium prices $p^{*}$ are regular.

\section*{Finitely many equilibria}
Theorem 1 Under the maintained hypotheses, every regular economy has finitely many equilibrium prices (up to normalization).

\begin{itemize}
  \item See Mas-Colell, Whinston, and Green (1995), Proposition 17.D.1.
  \item This implies that the Walrasian equilibrium prices are locally unique.
\end{itemize}

\section*{References}
Debreu, Gerard. 1970. "Economies with a Finite Set of Equilibria." Econometrica 38: 387-92. Mas-Colell, Andreu, Michael Whinston, and Jerry Green. 1995. Microeconomic Theory. New York, Oxford: Oxford University Press.


\end{document}